% Reader question: can consistent individual rankings yield a majority cycle?
% Claim: this profile gives a Condorcet cycle; cardinal sums use extra assumptions.
% Evidence: three explicit rankings, compatible utility bids and their totals.
% Distinguish: majority non-transitivity, Arrow hypotheses, cardinal comparability.
% Grammar: paired triangles on the upstream grid, not an unconditional escape.
% Acceptance: recount each 2-to-1 vote and all three cardinal sums.
\noindent\begin{minipage}{11.1cm}\centering
\pdfigurehue{pdcobalt}%
\begin{tikzpicture}[pd figure,x=1cm,y=1cm]
  % Paired triangles have the same geometry and native label size.
  \def\PX{5.65}\def\xA{2.55}\def\xB{1.15}\def\xC{4.00}
  \def\ytop{1.65}\def\ybase{-0.25}
  \node[pd title,anchor=west] at (.10,4.45) {ordinal ranking};
  \node[pd title,anchor=west] at (\PX,4.45) {cardinal utility};
  \foreach \i/\s in {0/{$A\succ B\succ C$},1/{$B\succ C\succ A$},2/{$C\succ A\succ B$}}{
    \node[pd label,anchor=west] at (.10,{3.90-\i*.45}) {agent \number\numexpr\i+1\relax:\ \s};}
  \foreach \i/\s in {0/{$A\,.9$\ $B\,.5$\ $C\,.3$},1/{$A\,.4$\ $B\,.9$\ $C\,.8$},2/{$A\,.6$\ $B\,.2$\ $C\,.9$}}{
    \node[pd label,anchor=west] at (\PX,{3.90-\i*.45}) {agent \number\numexpr\i+1\relax:\ \s};}
  \draw[pd hairline] (\PX,2.80) -- (10.80,2.80);
  \node[pd title,anchor=west] at (\PX,2.50) {$\Sigma$: $A\,1.9$\ $B\,1.6$\ $C\,2.0$};
  \node[pd automaton state] (La) at (\xA,\ytop) {$A$};
  \node[pd automaton state] (Lb) at (\xB,\ybase) {$B$};
  \node[pd automaton state] (Lc) at (\xC,\ybase) {$C$};
  \draw[pd breach arrow] (La) -- (Lb);
  \draw[pd breach arrow] (Lb) -- (Lc);
  \draw[pd breach arrow] (Lc) -- (La);
  \node[pd breach label] at (\xA,-1.15) {$A\succ B\succ C\succ A$: no winner};
  \node[pd automaton state] (Ra) at ({\PX+\xA},\ytop) {$A$};
  \node[pd automaton state] (Rb) at ({\PX+\xB},\ybase) {$B$};
  \node[pd automaton focus] (Rc) at ({\PX+\xC},\ybase) {$C$};
  \draw[pd focus arrow] (Rc) -- (Ra);
  \draw[pd focus arrow] (Ra) -- (Rb);
  \draw[pd focus arrow] (Rc) -- (Rb);
  \node[pd focus label] at ({\PX+\xA},-1.15) {$C\succ A\succ B$: total order};
\end{tikzpicture}
\par\smallskip
{\pdfiglabel\renewcommand{\arraystretch}{1.15}%
\begin{tabularx}{\linewidth}{@{}p{3.25cm}>{\raggedright\arraybackslash}X@{}}
\toprule
\textbf{Condition} & \textbf{Ordinal scope}\\
\midrule
U: unrestricted domain & All strict ranking profiles are allowed.\\
P: weak Pareto & Unanimous pairwise preferences are respected.\\
IIA: independence & A pair depends only on its individual comparisons.\\
D: dictatorship & Forced by U, P, IIA and complete transitivity.\\
\bottomrule
\end{tabularx}}
\par\smallskip
{\pdfiglabelitalic\raggedright Cycles are possible, not inevitable. Arrow assumes a complete,
transitive social order. Cardinal summation uses additional scale and comparability
assumptions; it gives no unconditional governance guarantee.\par}
\end{minipage}
